\chapter{Lösungen der Übungsaufgaben}

\structhead{Aufgabe 1.1.1.6:}
Die aktuelle Gesamtschuld $S$ ist die Summe der beiden Barwerte:
\[ S = \frac{3.000}{1 + 0{,}07 \cdot \frac{5}{12}} + \frac{2.000}{1 + 0{,}07 \cdot \frac{10}{12}} = 2.914{,}98 + 1.889{,}76 = 4.804{,}74 \]

\structhead{Aufgabe 1.1.3.5:}
Man berechnet zuerst den Betrag nach vier Jahren:
\[ K_4 = 2.000 \cdot 1{,}03^4 = 2.251{,}02 \]
Hierzu kommen noch einmal Zinsen in Höhe von $\dfrac{117}{365}$ von 3\,\%:
\[ Z = 2.251{,}02 \cdot 0{,}03 \cdot \frac{117}{365} \approx 21{,}65 \]
Insgesamt ergibt sich so ein Betrag von
\[ K_{4/117T} = 2.272{,}67\,€ . \]

\structhead{Aufgabe 1.1.3.7:}
Um die Angebote vergleichen zu können, muss man sie alle auf den gleichen Zeitpunkt beziehen. Nimmt man als diesen Zeitpunkt den Punkt 0 (das entspricht »heute«), so muss man die Barwerte der drei vorliegenden Angebote bestimmen. Dies führt hier zu folgenden Ergebnissen:
\begin{itemize}
  \item Barwert von Angebot A: 140.000\,€
  \item Barwert von Angebot B: $100.000 + \dfrac{60.000}{1{,}08^5} \approx 140.834{,}99\,€$
  \item Barwert von Angebot C: $\dfrac{175.000}{1{,}08^3} \approx 138.920{,}64\,€$
\end{itemize}
Angebot B ist also das günstigste für den Verkäufer. Hätte man alle Angebote auf einen anderen Zeitpunkt -- z.\,B. in 3 Jahren -- bezogen, so hätten sich selbstverständlich andere Werte ergeben (alle mit $1{,}08^3$ multipliziert), jedoch ist auch dann B das beste Angebot.

\structhead{Aufgabe 1.1.4.2:}
Ein nomineller Jahreszinssatz von 6\,\% führt mit vierteljährlicher Verzinsung zu einem effektiven Zinssatz von
\[ i_E = 1{,}015^4 - 1 = 0{,}0614 \text{ bzw. } 6{,}14\,\% . \]
Ein monatlicher relativer Zinssatz von 0{,}5\,\% ergibt jährlich
\[ i_E = 1{,}005^{12} - 1 = 0{,}0617 \text{ bzw. } 6{,}17\,\% . \]
Das letzte Angebot ist also das beste.

\structhead{Aufgabe 1.2.1.5:}
Durch die 60-malige Zahlung der Rate 300\,€ ergibt sich ein Rentenendwert
\[ R_{60} = 300 \cdot \frac{\left(1 + \frac{0{,}075}{12}\right)^{60} - 1}{\left(1 + \frac{0{,}075}{12}\right) - 1} \approx 21.758{,}13 . \]
Der Barwert beträgt
\[ R_0 = \frac{21.758{,}13}{1{,}075^5} \approx 15.155{,}81 . \]
Die Studentin L. kann also einen Kreditbetrag von 15.155{,}81\,€ aufnehmen.

\structhead{Aufgabe 1.2.1.8:}
Einsetzen der gegebenen Werte in die Formel
\[ q^n - \frac{R_n^{\text{nach}}}{r} \cdot q + \frac{R_n^{\text{nach}}}{r} - 1 = 0 \]
ergibt hier die Gleichung
\[ q^3 - 3{,}10623\,q + 2{,}10623 = 0 . \]
Mit Hilfe eines geeigneten Taschenrechners oder eines Computeralgebrasystems erhält man drei Lösungen: $q_1 \approx -2{,}035$, $q_2 = 1$ und $q_3 \approx 1{,}035$. Nur die dritte Lösung kommt in Frage, der gesuchte Zinssatz beträgt folglich 3{,}5\,\%.

\structhead{Aufgabe 1.2.1.11:}
Als Endwert der Einzahlungen ergibt sich:
\[ R_{20}^{\text{vor}} = 2.000 \cdot 1{,}06 \cdot \frac{1{,}06^{20} - 1}{1{,}06 - 1} = 77.985{,}45\,€ \]
Dies ist gleichzeitig der Barwert zu Beginn der Auszahlungsphase. Für die Auszahlungsrate erhält man:
\[ r = 77.985{,}45 \cdot \frac{1{,}05^{15} \cdot 0{,}05}{1{,}05^{15} - 1} = 7.513{,}30\,€ \]
Die Person wird also fünfzehn Jahre lang jeweils am Jahresende 7.513{,}30\,€ ausgezahlt bekommen, dann ist das Guthaben aufgebraucht.

\structhead{Aufgabe 1.2.2.4:}
Zunächst wird berechnet, welcher Betrag jährlich (nachschüssig) einer vorschüssigen monatlichen Rate von 200\,€ entspricht:
\[ \bar{r}_E = 200 \left(12 + \tfrac{0{,}06}{2}(12 + 1)\right) = 2.478 \]
Nun ist der Rentenendwert der über 40 Jahre nachschüssig gezahlten Rente in Höhe von $r = 2.478$ zu bestimmen:
\[ R_{40}^{\text{nach}} = 2.478 \cdot \frac{1{,}06^{40} - 1}{1{,}06 - 1} = 383.500{,}15 \]
Es stehen also 383.500{,}15\,€ zur Verfügung.

\structhead{Aufgabe 1.3.1.3:}
Der Ansatz
\[ 50.000 \cdot 1{,}08^{15} = A \cdot \frac{1{,}08^{15} - 1}{1{,}08 - 1} \]
führt zu
\[ A = 5.841{,}48 . \]
Am Jahresende müssen also jeweils 5.841{,}48\,€ gezahlt werden.

\structhead{Aufgabe 1.3.1.5:}
\begin{center}
  \centering
  \begin{tabular}{rrrrr}
    \hline
    Jahr & Restschuld (zu Beginn des Jahres) & Zinsen $Z_j$ & Tilgung $T_j$ & Annuität $A_j$\\
    \hline
    1 & 30.000 & 2.250 & 6.000 & 8.250\\
    2 & 24.000 & 1.800 & 6.000 & 7.800\\
    3 & 18.000 & 1.350 & 6.000 & 7.350\\
    4 & 12.000 & 900 & 6.000 & 6.900\\
    5 & 6.000 & 450 & 6.000 & 6.450\\
    6 & 0 & & & \\
    \hline
  \end{tabular}
\end{center}

\structhead{Aufgabe 1.4.1.3:}
\begin{center}
  \begin{tabular}{rrrrrr}
    \hline
    Periode $k$ & Einnahmen $E_k$ & Ausgaben $A_k$ & Überschüsse $E_k - A_k$ & $\dfrac{1}{1{,}07^k}$ & Barwerte der Überschüsse\\
    \hline
    0 & 0 & 435.000 & $-435.000$ & 1 & $-435.000{,}00$\,€\\
    1 & 150.000 & 45.000 & 105.000 & 0{,}934579 & 98.130{,}84\,€\\
    2 & 180.000 & 60.000 & 120.000 & 0{,}873439 & 104.812{,}65\,€\\
    3 & 210.000 & 80.000 & 130.000 & 0{,}816298 & 106.118{,}72\,€\\
    4 & 190.000 & 70.000 & 120.000 & 0{,}762895 & 91.547{,}43\,€\\
    5 & 170.000 & 65.000 & 105.000 & 0{,}712986 & 74.863{,}55\,€\\
    \hline
    \multicolumn{5}{r}{Kapitalwert $C$:} & 40.473{,}19\,€\\
    \hline
  \end{tabular}
\end{center}
Die Investition ist wegen $C > 0$ als vorteilhaft einzuschätzen.

\structhead{Aufgabe 1.4.2.3:}
Die Bedingung $C = 0$ führt in diesem Fall nach Multiplikation mit $(-q^5)$ und Division durch 1.000 zu der folgenden Gleichung:
\[ 435\,q^5 - 105\,q^4 - 120\,q^3 - 130\,q^2 - 120\,q - 105 = 0 \]
Man braucht nun einen guten Taschenrechner (den man auch bedienen kann) oder ein Computeralgebrasystem, um zu der einzigen reellen Lösung zu kommen:
\[ q \approx 1{,}104 \]
Der interne Zinsfuß liegt also bei ca. 10{,}4\,\%.
